\subsection{Optimising the harvest policy}\label{sec:yield}
Finally we let a genetic algorithm design the prey-harvest policy: the genome holds the controller type
(effort, quota or feedback) and its parameters, and fitness is the mean prey yield per 1000 cells in the
second half of a 1500-tick run on a $100^2$ grid (3 training seeds, population 32, \Ygens{} generations,
\Yevals{} simulations). No stability or coexistence term was included. The search converged on
\YmodeName{} with $E=\Yeffort$, a training fitness of \Ytrain{} prey per tick per 1000 cells. Among
policies that kept both species alive on all training seeds, the best effort, quota and feedback designs
reached \YbestEff, \YbestQuo{} and \YbestFb{}, respectively. Feedback was abandoned early: only 11 of
632 genomes used it, and 9 of those had set points above the prey level of a $100^2$ world, so the
controller was idle. The GA's verdict on feedback is therefore uninformative, not negative.

Re-validated on 10 fresh seeds, the same genome scored only \Yval{}: predators survived in \Ycoex\% of
seeds (mean \YmeanPred{} predators). The evolved effort sits just below the point where predators are lost
in E6a, the edge of the coexistence region, where the outcome depends on chance. This is the expected
behaviour of an optimiser given a yield-only objective with few seeds: it climbs to the boundary of the
feasible region and \emph{overfits} the training noise. Robust management requires putting the constraint
(coexistence, variability) into the objective and evaluating it over many stochastic realisations---the
same lesson as the specification gaming of Sec.~\ref{sec:opt}, now with an economic objective.
